\paragraph{Contextual adaptation from glyph images.}
The starting reader comprised a fitted glyph-image metric, a continuous visual projection and a 12-layer Chinese-pretrained question-answering encoder. Its word embedding table was absent. For rank four maps $A$ and $B$, eight attention projections computed $Wx+B(Ax)$, with $B$ initially zero. We trained the query and value residuals in the last four layers and the boundary readout. A candidate reply $g$ was a contiguous strip of at most eight source cells. Its score pooled all occurrences by log-sum-exp of their start and end logits. Overlapping contextual windows covered longer pages. Training alternated 600 classical batches and 600 batches of paired human novel questions. Both arms had the same novel targets. The classical arm used the annotated person in the joint page and two-name objectives, while the control averaged the objectives for the two people. Inference selected the highest-scoring image among all candidate strips. The source of each of 14 further classical excerpts was verified in the \textit{Discourses of the States} before fitting. We computed paired percentile intervals by resampling original excerpts 10,000 times, keeping both pages and questions together. The inherited encoder's original pretraining corpus could include these works. Data preparation, base pretraining and evaluation were excluded from adapter fit time and energy.
